\chapter{Mongoose Schemas and Models}

\section{Schema vs Model}

\whatwhyhow{
A \textbf{schema} is the blueprint for a document's shape and validation rules. A \textbf{model} is the compiled interface Mongoose builds from that schema to query a collection.
}{
Separating the two lets you define validation once while getting a full query API (\texttt{find}, \texttt{create}, \texttt{findById}, ...) for free.
}{
Call \texttt{new mongoose.Schema(\{...\})} to define fields, then \texttt{mongoose.model("Task", taskSchema)} to compile it (Mongoose pluralizes/lowercases the collection name automatically).
}

\section{Task Schema}

\begin{lstlisting}[language=JavaScript]
// models/Task.js
const mongoose = require("mongoose");

const taskSchema = new mongoose.Schema(
  {
    title: { type: String, required: [true, "Title is required"], trim: true },
    description: { type: String, default: "" },
    status: { type: String, enum: ["todo", "progress", "completed"], default: "todo" },
    priority: { type: String, enum: ["low", "medium", "high"], default: "medium" },
    dueDate: { type: Date },
    user: { type: mongoose.Schema.Types.ObjectId, ref: "User", required: true },
  },
  { timestamps: true }
);

module.exports = mongoose.model("Task", taskSchema);
\end{lstlisting}

\section{Field Options Explained}

\begin{table}[h]
\centering
\begin{tabularx}{\textwidth}{l X}
\toprule
\textbf{Option} & \textbf{Meaning} \\
\midrule
\texttt{type}      & Data type cast/validated against \\
\texttt{required}   & Rejects save if missing \\
\texttt{enum}       & Restricts to fixed allowed values \\
\texttt{default}    & Value used if not supplied \\
\texttt{ref}        & Model an ObjectId points to, enables \texttt{.populate()} \\
\texttt{trim}       & Strips whitespace from Strings \\
\texttt{timestamps} & Schema option adding \texttt{createdAt}/\texttt{updatedAt} \\
\bottomrule
\end{tabularx}
\end{table}

\section{User Schema (for later chapters)}

\begin{lstlisting}[language=JavaScript]
// models/User.js
const userSchema = new mongoose.Schema(
  {
    name: { type: String, required: [true, "Name is required"], trim: true },
    email: { type: String, required: true, unique: true, lowercase: true, trim: true },
    password: { type: String, required: true, minlength: 6, select: false },
    avatar: { type: String, default: "" },
    role: { type: String, enum: ["user", "admin"], default: "user" },
  },
  { timestamps: true }
);

module.exports = mongoose.model("User", userSchema);
\end{lstlisting}

\begin{notebox}[NOTE]
\texttt{unique: true} on \texttt{email} creates a DB-level unique index, not just a validation rule.
\end{notebox}

\begin{tipbox}[TIP]
\texttt{select: false} on \texttt{password} means queries never return the hash by default; opt in with \texttt{.select("+password")} during login.
\end{tipbox}
